% FCS Research Article working draft. Generated results: experiments/fcs/.
\documentclass[research]{fcs}
\usepackage[T1]{fontenc}
\usepackage{booktabs}
\usepackage{placeins}
\usepackage{hyperref}
\setcounter{dbltopnumber}{4}
\renewcommand{\dbltopfraction}{0.85}
\renewcommand{\textfraction}{0.1}
\newcommand{\RefBase}{86.03}
\newcommand{\RefBaseSD}{1.40}
\newcommand{\RefNovel}{66.90}
\newcommand{\RefNovelSD}{4.23}
\newcommand{\RefAll}{76.46}
\newcommand{\RefAllSD}{1.43}
\newcommand{\RefHarmonic}{75.19}
\newcommand{\RefHarmonicSD}{2.14}
\newcommand{\DynamicBase}{85.36}
\newcommand{\DynamicBaseSD}{0.66}
\newcommand{\DynamicNovel}{70.83}
\newcommand{\DynamicNovelSD}{4.33}
\newcommand{\DynamicAll}{78.09}
\newcommand{\DynamicAllSD}{2.02}
\newcommand{\DynamicHarmonic}{77.36}
\newcommand{\DynamicHarmonicSD}{2.48}
\newcommand{\SimpleBase}{82.47}
\newcommand{\SimpleBaseSD}{0.87}
\newcommand{\SimpleNovel}{77.53}
\newcommand{\SimpleNovelSD}{0.47}
\newcommand{\SimpleAll}{80.00}
\newcommand{\SimpleAllSD}{0.56}
\newcommand{\SimpleHarmonic}{79.92}
\newcommand{\SimpleHarmonicSD}{0.55}
\newcommand{\TextBase}{78.32}
\newcommand{\TextBaseSD}{0.25}
\newcommand{\TextNovel}{76.86}
\newcommand{\TextNovelSD}{1.97}
\newcommand{\TextAll}{77.59}
\newcommand{\TextAllSD}{0.86}
\newcommand{\TextHarmonic}{77.57}
\newcommand{\TextHarmonicSD}{0.88}
\newcommand{\RefErrorFlip}{1.84}
\newcommand{\RefErrorFlipSD}{1.54}
\newcommand{\RefChurn}{3.63}
\newcommand{\RefChurnSD}{2.40}
\newcommand{\LatencyFifty}{1.76}
\newcommand{\LatencySixty}{0.75}
\newcommand{\LatencyEighty}{0.78}
\newcommand{\PetRefBase}{97.18}
\newcommand{\PetRefBaseSD}{2.25}
\newcommand{\PetRefNovel}{87.79}
\newcommand{\PetRefNovelSD}{6.18}
\newcommand{\PetRefAll}{92.49}
\newcommand{\PetRefAllSD}{2.83}
\newcommand{\PetRefHarmonic}{92.14}
\newcommand{\PetRefHarmonicSD}{3.23}
\newcommand{\PetDynamicBase}{96.85}
\newcommand{\PetDynamicBaseSD}{1.50}
\newcommand{\PetDynamicNovel}{86.49}
\newcommand{\PetDynamicNovelSD}{7.64}
\newcommand{\PetDynamicAll}{91.67}
\newcommand{\PetDynamicAllSD}{4.56}
\newcommand{\PetDynamicHarmonic}{91.29}
\newcommand{\PetDynamicHarmonicSD}{4.99}
\newcommand{\PetSimpleBase}{94.77}
\newcommand{\PetSimpleBaseSD}{2.91}
\newcommand{\PetSimpleNovel}{93.91}
\newcommand{\PetSimpleNovelSD}{1.55}
\newcommand{\PetSimpleAll}{94.34}
\newcommand{\PetSimpleAllSD}{1.43}
\newcommand{\PetSimpleHarmonic}{94.31}
\newcommand{\PetSimpleHarmonicSD}{1.44}
\newcommand{\PetTextBase}{92.87}
\newcommand{\PetTextBaseSD}{3.37}
\newcommand{\PetTextNovel}{94.05}
\newcommand{\PetTextNovelSD}{1.35}
\newcommand{\PetTextAll}{93.46}
\newcommand{\PetTextAllSD}{1.59}
\newcommand{\PetTextHarmonic}{93.43}
\newcommand{\PetTextHarmonicSD}{1.63}

\title{Reference-anchored candidate tables for class expansion in long-tailed recognition}
\author[1]{Junyi Zhang}
\address[1]{Shandong University, China}
\corremail{cheerssom0876@gmail.com}
\fcssetup{received={not submitted},accepted={not assigned},
  copyright={Working draft, October 21, 2026},
  publish-year=2026,publish-month={October},article-id={Draft}}
% Suppress the publisher's example DOI and issue metadata in a working draft.
\ExplSyntaxOn
\tl_set:Nn \l__fcs_doi_url_tl { FCS~Research~Article~---~Working~Draft }
\tl_set:Nn \l__fcs_short_detailed_page_tl { FCS~submission~draft }
\tl_set:Nn \l__fcs_long_detailed_page_tl { FCS~submission~draft~\textbar~\thepage }
\ExplSyntaxOff

\begin{abstract}
Expandable recognition requires a consistent score system as new classes join an imbalanced base vocabulary. Combining semantic similarities, learned logits, and visual prototypes can disrupt this consistency when normalization depends on the active candidates. We introduce reference-anchored candidate tables, which represent each class by its available evidence and normalize each source against fixed base-class references. A new class enters through its name and support images, without parameter updates or gate recalibration. We prove that append-only expansion preserves existing scores and their relative ordering, while allowing new candidates to compete for the final prediction. Frozen-CLIP evaluations on CIFAR-100 and Oxford-IIIT Pets use three class partitions per dataset, disjoint support images, and base-only calibration. At five shots, the table attains \RefBase\% base and \RefNovel\% novel accuracy on CIFAR-100, and \PetRefBase\%/\PetRefNovel\% on Pets. Text--prototype fusion attains \SimpleBase\%/\SimpleNovel\% and \PetSimpleBase\%/\PetSimpleNovel\%, respectively, placing the table at a retention-oriented operating point. Vocabulary tests measure zero base-score drift on both datasets; five-session expansion gives identical final predictions across two arrival orders in each partition. Calibration controls and evidence ablations connect this behavior to the reference frame and the available evidence. The resulting interface combines trained discrimination with sparse novel-class support in a stable score system.
\end{abstract}
\keywords{Long-tailed recognition; Class expansion; Vision--language models; Evidence fusion; Score stability; Few-shot classification}

\begin{document}

\section{Introduction}
Recognition systems acquire new categories after deployment while retaining classes learned from imbalanced data. Long-tailed recognition addresses the unequal supervision available to common and rare classes~\cite{liu2019large,cao2019ldam,cui2019classbalanced}. Expansion adds a second requirement: a classifier must combine established discrimination with arriving evidence, often a class name and a handful of images.

Pre-trained vision--language models make class names immediately useful for recognition. CLIP maps images and descriptions into a shared space~\cite{radford2021clip}; prompt adaptation and lightweight fine-tuning specialize this representation~\cite{zhou2022coop,zhou2022cocoop,shi2024lift}; feature caches and prototypes incorporate labeled support~\cite{zhang2022tipadapter,snell2017prototypical,wang2019simpleshot}. An expandable recognizer can therefore combine trained weights for base classes with semantic anchors and visual prototypes for arriving classes.

The score system connecting these sources must remain consistent as the vocabulary grows. Learned logits, text similarities, and prototype similarities have different scales. Standardizing each branch over its active candidates makes those scales depend on the vocabulary: adding a category changes the normalization of existing text scores, even with frozen parameters. In multi-branch fusion, this rescaling can reverse existing-class rankings.

We introduce \emph{reference-anchored candidate tables} to decouple evidence normalization from vocabulary growth. Each row stores the evidence available for a class, and each branch uses a fixed bank of base-class references. Appending a row adds a competitor while preserving every existing score in real arithmetic. The table combines heterogeneous supervision within a shared reference frame.

Class-incremental learning studies the retention of earlier discrimination~\cite{rebuffi2017icarl,tao2020fewshot}, and prototype construction with pre-trained representations enables efficient class addition~\cite{zhou2024simplecil,mcdonnell2023ranpac}. Our contribution concerns the fusion stage: fixed references preserve the score system across classes with different evidence types. Existing-class rankings remain stable, and a prediction changes precisely when an appended candidate wins.

Across CIFAR-100 and Oxford-IIIT Pets, a frozen CLIP encoder supplies features to a base-trained classifier and novel support prototypes. Base-only pseudo-expansion episodes select the fusion coefficients. At five shots, base/novel accuracy is \RefBase\%/\RefNovel\% on CIFAR and \PetRefBase\%/\PetRefNovel\% on Pets. Vocabulary and sequential-arrival tests verify score stability; offset-calibrated controls and retention-targeted selection characterize the accompanying recognition trade-off.

The contributions are an availability-masked candidate table, fixed-reference fusion with a score-preservation proof, and a base-only calibrated evaluation of recognition, sequential expansion, and readout cost.

\section{Related work}
\subsection{Long-tailed recognition and calibration}
Long-tailed recognition addresses the bias induced by unequal class frequencies. Margin adjustment allocates different classification margins to frequent and rare classes~\cite{cao2019ldam}; class-balanced weighting uses effective sample counts~\cite{cui2019classbalanced}; and prior-aware objectives modify the relationship between logits and class probabilities~\cite{ren2020bsm,menon2021logitadj}. Representation--classifier decoupling provides a complementary route: a representation is learned first and a balanced classifier is then fitted to that representation~\cite{kang2020decoupling}. These developments motivate our use of a balanced base classifier as one source of evidence.

Foundation models change the representation-learning budget and the role of task-specific adaptation. LIFT studies lightweight adaptation of pre-trained models for long-tail learning~\cite{shi2024lift}. Low-rank updates provide an efficient mechanism for adapting large parameter matrices~\cite{hu2022lora}, including transformer backbones~\cite{dosovitskiy2021vit}. Our candidate interface can accept weights obtained through these routes. The expansion experiments here fix the encoder throughout and isolate the effect of the readout.

Long-tailed calibration studies confidence under class imbalance~\cite{zhong2021mislas}, while robust long-tailed learning addresses its interaction with label noise~\cite{wei2026robust}. We study the reference population used to standardize evidence as candidate classes change. Class-balanced validation and separate base/novel metrics expose how the resulting score system allocates discrimination.

\subsection{Semantic and visual evidence for few-shot recognition}
CLIP constructs classifiers from natural-language supervision and supports classification by changing the text candidates~\cite{radford2021clip}. CoOp learns prompt context for target tasks, while CoCoOp introduces conditional prompting to improve transfer across categories~\cite{zhou2022coop,zhou2022cocoop}. These methods concern the construction or adaptation of semantic evidence. Our table keeps semantic anchors available to every class and combines them with evidence that may exist for only a subset of the vocabulary.

Prototype methods summarize class examples in an embedding space~\cite{snell2017prototypical}. SimpleShot studies feature normalization for nearest-neighbor classification~\cite{wang2019simpleshot}, and Tip-Adapter augments CLIP with a support-feature cache~\cite{zhang2022tipadapter}. Few-shot transfer under distribution shift further motivates separating shared and task-specific knowledge~\cite{zhu2023combat}. Our comparisons use a shared encoder and novel support budget to examine evidence availability and its normalization reference.

\subsection{Incremental recognition with pre-trained models}
Class-incremental learning evaluates a model on the union of classes observed across successive training stages. Exemplar rehearsal and representation constraints are established tools for retaining earlier knowledge~\cite{rebuffi2017icarl}. Few-shot class-incremental learning adds a data scarcity condition to the arriving classes~\cite{tao2020fewshot}. Recent FCS work studies temporal consistency with unlabeled memory replay~\cite{wang2025memory} and the geometry of inter-task and intra-task discrimination~\cite{tang2027global}. These approaches motivate explicit measurement of existing-class behavior.

SimpleCIL examines prototype classifiers built from pre-trained representations~\cite{zhou2024simplecil}; RanPAC combines random projections with accumulated class statistics~\cite{mcdonnell2023ranpac}. Reference-anchored fusion extends this setting to semantic, trained, and support-derived evidence with different availability patterns. Calibrated seen/unseen recognition~\cite{chao2016calibrated} provides a related perspective on the base--novel operating trade-off.

TV100 studies categories unavailable during CLIP's original pre-training period~\cite{zhou2024tv100}. Here, novelty is defined relative to task-specific classifier fitting and calibration; CLIP retains its web-scale pre-training exposure.

\section{Expandable recognition with heterogeneous evidence}
\subsection{Problem definition}
Let $\mathcal B$ denote base classes with a labeled training set $\mathcal D_B$, and let $\mathcal N$ denote novel classes supplied at expansion time. The active vocabulary is initially $\mathcal C_0=\mathcal B$ and subsequently becomes $\mathcal C=\mathcal B\cup\mathcal N$. An image encoder $f$ maps an input image to a normalized vector $v=f(x)\in\mathbb R^D$. A text encoder supplies a normalized class anchor $t_c$. For each $c\in\mathcal N$, the support set $\mathcal S_c$ contains $k$ labeled examples. Inference is over all active classes, without an oracle indicating whether an image belongs to a base or novel class.

Base training can be imbalanced. If $n_c$ is the number of training images for class $c$, the classifier may use inverse-frequency weights or another balanced fitting procedure. After this stage, the encoder, base classifier, references, and fusion coefficients are fixed. At arrival, novel class information is restricted to its name and $\mathcal S_c$. Neither previous training images nor additional labeled novel validation images are required to construct its row.

Each row has an availability mask $m_{ec}\in\{0,1\}$ over evidence sources $e$. We use semantic evidence $T$, learned discriminative evidence $L$, and visual prototype evidence $P$. The experimental configuration is
\begin{equation}
 (m_{Tc},m_{Lc},m_{Pc})=
 \begin{cases}(1,1,0),&c\in\mathcal B,\\(1,0,1),&c\in\mathcal N.\end{cases}
 \label{eq:availability}
\end{equation}
Base prototypes form the visual branch's normalization references. The availability mask determines which evidence contributes to a candidate; the reference bank determines its scale.

\begin{figure*}[t]
\centering
\includegraphics[width=.98\textwidth]{fig1_framework_fcs.pdf}
\caption{Reference-anchored candidate-table architecture. A shared frozen CLIP feature supplies semantic similarity and the base-trained readout. The candidate table stores text anchors for all classes, learned weights for base classes, and support prototypes for arriving classes. The inset shows the three evidence branches normalized against their respective fixed base reference sets before availability-masked fusion; $z_e=\widehat q_e$ denotes normalized evidence. A novel row is appended without updating the encoder, existing weights, references, or gates. Orange identifies base-fitted evidence and fusion; blue hatching identifies pre-trained or cached evidence.}
\label{fig:framework}
\end{figure*}

\subsection{Three notions of preservation}
Preserving a parameter vector $\theta$ means that $\theta$ receives no update during expansion. Preserving a score additionally requires
\begin{equation}
 s_{\mathcal C_0}(x,c)=s_{\mathcal C}(x,c),\quad c\in\mathcal C_0.
 \label{eq:scorepreservation}
\end{equation}
Preserving a prediction requires the maximizing class over the larger vocabulary to remain the same. The third property is stronger than the second: a newly admitted class can have a larger score than the previous winner. We seek score preservation while allowing that competition.

The evaluation separates changes among existing classes from selection of an appended class. Under fixed evidence, reference anchoring preserves existing-class competition and admits new alternatives through their own scores.

\section{Reference-anchored candidate tables}
\subsection{Evidence construction}
The semantic score is cosine similarity between a normalized image feature and a class text anchor,
\begin{equation}
 q_T(x,c)=v^\top t_c.
\end{equation}
We average normalized text embeddings of eight fixed prompts and normalize the result. The same prompts construct anchors for base and arriving classes.

The learned score is the output of a base classifier,
\begin{equation}
 q_L(x,c)=v^\top w_c+b_c,\quad c\in\mathcal B.
\end{equation}
For the frozen-encoder experiments, the weights are fitted using class-balanced ridge regression. With augmented image features $\widetilde v=[v;1]$, one-hot targets $Y$, and diagonal sample-weight matrix $\Omega$, the solution is
\begin{equation}
 W=(X^\top\Omega X+\lambda R)^{-1}X^\top\Omega Y,
 \label{eq:ridge}
\end{equation}
where $R=\operatorname{diag}(1,\ldots,1,0)$ leaves the intercept unregularized. A sample from class $c$ has weight $1/n_c$, giving each base class equal total fitting weight. We solve the corresponding linear system.

The visual prototype of a novel class is
\begin{equation}
 p_c=\frac{\sum_{x\in\mathcal S_c}f(x)}{\|\sum_{x\in\mathcal S_c}f(x)\|_2},
 \qquad q_P(x,c)=v^\top p_c.
 \label{eq:prototype}
\end{equation}
The same averaging rule applied to base training features produces a reference bank $\{p_b:b\in\mathcal B\}$. The novel prototype stores the information acquired at class arrival; the base bank fixes the coordinate system in which its score is interpreted.

\subsection{Why a changing normalization population matters}
Candidate-wise standardization computes the mean and standard deviation over the current valid candidates of a branch. For text evidence, this gives
\begin{equation}
 z_T^{\mathcal C}(x,c)=\frac{q_T(x,c)-\mu_T(x;\mathcal C)}{\max(\sigma_T(x;\mathcal C),\epsilon)}.
 \label{eq:dynamic}
\end{equation}
Adding candidates changes $\mu_T$ and $\sigma_T$. For a text-only classifier, a common positive affine transformation preserves the ordering of the old classes. With heterogeneous fusion, however, the learned branch and text branch are rescaled differently. The relative strength of their votes can change.

A two-class example illustrates the effect. Suppose the old text scores are $(0,2)$, giving standardized scores $(-1,1)$, while the learned branch contributes $(0.25,-0.25)$. The second old class has the larger fused score. Appending a text score of $100$ changes the two old standardized scores to approximately $(-0.728,-0.685)$. Their text separation is now about $0.043$, and the unchanged learned contribution reverses the old-class ordering. No classifier weight has changed. The example isolates a normalization effect independently of whether the new class ultimately wins the expanded decision.

\subsection{Fixed references for each evidence source}
Let $\mathcal R_e$ be a fixed bank of reference rows for source $e$. For every query, compute
\begin{align}
 \mu_e(x)&=\frac{1}{|\mathcal R_e|}\sum_{r\in\mathcal R_e}q_e(x,r),\\
 \sigma_e^2(x)&=\frac{1}{|\mathcal R_e|}\sum_{r\in\mathcal R_e}(q_e(x,r)-\mu_e(x))^2.
 \label{eq:reference}
\end{align}
The normalized candidate score is
\begin{equation}
 \widehat q_e(x,c)=\frac{q_e(x,c)-\mu_e(x)}{\max(\sigma_e(x),\epsilon)}.
 \label{eq:anchored}
\end{equation}
The reference statistics remain query-dependent, which adapts the scale to the current image, but are independent of which novel candidates are active. We use base text anchors for $\mathcal R_T$, base classifier columns for $\mathcal R_L$, and base visual means for $\mathcal R_P$. Each bank contains 40 reference classes on CIFAR-100 and 15 on Pets. A zero-variance reference bank uses the same fixed denominator floor $\epsilon=10^{-6}$.

The candidate score combines available normalized evidence:
\begin{equation}
 s(x,c)=\sum_{e\in\{T,L,P\}}m_{ec}\alpha_e\widehat q_e(x,c)+\delta_c.
 \label{eq:fusion}
\end{equation}
We set $\alpha_T=1$, write $a=\alpha_L$ and $u=\alpha_P$, and use $\delta_c=d\,\mathbf 1[c\in\mathcal N]$. The resulting implementation is
\begin{align}
 s(x,b)&=\widehat q_T(x,b)+a\widehat q_L(x,b),\quad b\in\mathcal B,\\
 s(x,n)&=\widehat q_T(x,n)+u\widehat q_P(x,n)+d,\quad n\in\mathcal N.
 \label{eq:base-novel}
\end{align}
The offset $d$ calibrates competition between base and novel classes. The availability mask excludes absent learned evidence from the novel score and from reference statistics.

\subsection{Score preservation under append-only expansion}
\begin{proposition}
Fix an input $x$, the evidence of every existing row, the encoder, the reference banks, the denominator floor, and the coefficients in Eq.~\eqref{eq:fusion}. Appending new rows leaves the score of every existing row unchanged. Hence, every pairwise score difference and the maximizing row restricted to the existing vocabulary are unchanged, using a fixed tie-breaking rule.
\label{prop:invariance}
\end{proposition}
\begin{proof}
For an existing candidate $c$, each raw score $q_e(x,c)$ and mask $m_{ec}$ is fixed. The means and variances in Eq.~\eqref{eq:reference} are functions of $x$ and $\mathcal R_e$, and neither depends on the appended rows. Therefore each normalized score in Eq.~\eqref{eq:anchored} is unchanged. The coefficients and offset of $c$ are fixed, so their sum is unchanged. Subtracting the unchanged scores of any two existing rows gives an unchanged difference. A fixed tie-breaking rule preserves the maximizing existing row.
\end{proof}

The proposition covers both base rows and previously appended novel rows with fixed prototypes. The implementation computes the fixed base text block separately from novel columns to keep its numerical kernel independent of the appended width.

\begin{corollary}
If every existing row has a preserved score, a changed final prediction after expansion must select an appended row. An example whose previous prediction was incorrect cannot become correct after adding classes when its true class belongs to the existing vocabulary.
\label{cor:competition}
\end{corollary}
\begin{proof}
The best existing row is unchanged by Proposition~\ref{prop:invariance}. The expanded maximizer is either this row or a newly appended row. For an example labeled by an existing class, selecting a new row cannot correct a previous error, and keeping the existing maximizer leaves the previous error unchanged.
\end{proof}

\subsection{Base-only calibration}
Fusion calibration must respect the distinction between base supervision and novel support. We construct a pseudo-expansion episode by splitting the base classes into approximately equal pseudo-base and pseudo-novel groups: 20/20 on CIFAR-100 and 7/8 on Pets. A separate classifier is fitted using the pseudo-base classes alone. The pseudo-novel rows contain names and $k$ support examples from those base classes, while held-out base validation images supply query labels. This episode simulates missing learned evidence without consuming any actual novel-class validation labels.

For each support budget, candidate coefficients are evaluated on this episode. The selection objective is the harmonic mean of macro accuracy on the two groups,
\begin{equation}
 H=\frac{2A_BA_N}{A_B+A_N},
 \label{eq:harmonic}
\end{equation}
with $H=0$ if both terms vanish. The grid is $a,u\in\{0,0.25,0.5,1,2\}$ and $d\in\{-2,-1.5,\ldots,2\}$. Ties follow a fixed grid order. The full base classifier is then combined with actual novel rows using the selected coefficients. Neither actual novel validation queries nor test labels enter this selection rule.

This procedure uses one pseudo-class split per true partition and the same support-index seed as the evaluated support condition. The pseudo split seed is $100+p$ for true-partition seed $p$. Coefficients are chosen once for a support budget and stay fixed when the vocabulary-size stress test appends successive subsets of novel classes.

\subsection{Storage and computation}
The frozen implementation uses one image feature per input. Class text anchors, base classifier weights, visual references, and novel prototypes are cached. For feature dimension $D$, the base reference banks require $O(|\mathcal B|D)$ storage; each novel row requires $O(D)$ additional storage for its prototype and text anchor. Once its support images have been encoded, constructing a novel prototype requires $O(kD)$ arithmetic and no optimization state.

Scoring $K$ candidates requires $O(KD)$ arithmetic in addition to the fixed reference calculations. Image encoding is shared across the vocabulary, and existing row scores can be reused when appending candidates for the same query.

\section{Experiments}
\subsection{Data, access conditions, and implementation}
We evaluate expansion on CIFAR-100~\cite{krizhevsky2009cifar} with long-tailed base supervision. The training class counts follow an exponential schedule with imbalance parameter $\rho=100$ and nominal maximum 500 images. Five images per class are reserved before constructing the long-tailed pool, so the largest available pool contains 495 images. The remaining long-tailed pool contains 10,894 images, split into 10,487 fitting images and 407 held-out validation images. Validation reserves up to five examples per class, with smaller allocations for the rarest classes.

For seeds $\{0,1,2\}$, stratified sampling within head, medium, and tail frequency buckets yields distinct partitions of 40 base, 40 novel, and 20 unused classes. Proportional quotas use largest-remainder rounding after shuffling class membership. Evaluation uses the 80 base and novel classes.

Oxford-IIIT Pets~\cite{parkhi2012cats} provides a complementary fine-grained task with 37 cat and dog breeds. We use its official trainval/test split and reserve five training images per breed as support. From the remaining training images, we construct an exponential long-tailed pool with $\rho=10$ and maximum class count 90, giving 1,168 fitting and 146 validation images. Each true partition contains 15 base, 15 novel, and seven unused breeds, sampled with the same allocation rule. The official test split contains 3,669 images. Unlike CIFAR, its images have varied spatial resolutions; RGB images are resized to $224\times224$ by PIL bicubic interpolation before the same CLIP normalization. These features are freshly extracted for this study.

All expansion readouts use frozen CLIP ViT-L/14 features from a single $224\times224$ bicubic-resized image with CLIP channel normalization. Fitting, validation, and support pools are disjoint. Appendix~\ref{app:reproduction} gives feature-cache and image-pool verification details.

Base ridge regularization is selected from $\{0.01,0.1,1,10\}$ by base validation macro accuracy; the pseudo-base classifier uses the same rule. Novel support budgets are $k\in\{1,2,5\}$. One- and two-shot subsets use seeds $\{7,8,9\}$; five shots use the entire five-image reserve. We average support subsets within each partition and report the mean and sample standard deviation across the three distinct partitions.

Classifier fitting, regularization selection, and gate calibration use base classes exclusively. Five-shot CIFAR expansion consumes 200 support images for 40 novel classes. Its balanced test set supplies 4,000 base and 4,000 novel images for the expanded metrics.

\subsection{Readouts and evaluation measures}
All compared readouts use the same CLIP image embeddings, text anchors, class partitions, and novel support images. The text baseline uses cosine similarity to class names. The visual prototype baseline uses means of base fitting examples and means of novel support examples. Text--prototype fusion adds the two cosine branches with a mixing coefficient selected on the same pseudo-novel calibration episode; its grid is $\{0,0.25,0.5,1,2,4\}$, including pure text.

Tip-Adapter uses a balanced $k$-image cache per candidate and scores $100q_T+\alpha q_{\mathrm{cache}}$, with affinity $\exp[\beta(v^\top f(x_i)-1)]$. The same base-only episode selects $\alpha\in\{0,0.5,1,2,4\}$ and $\beta\in\{1,2,5,10\}$. Cache images come from the reserved pool for both groups; base-trained readouts use the full base fitting set.

We also compare raw evidence fusion and candidate-normalized fusion. Raw fusion uses unstandardized text, learned, and prototype scores with the same coefficient grid as the proposed table. Candidate-normalized fusion standardizes text over all active classes, learned scores over the base classes, and prototypes over active novel classes. It shares the base classifier and calibration protocol with reference-anchored fusion. This comparison changes the normalization population while holding the principal evidence sources fixed.

We report macro base accuracy $A_B$, macro novel accuracy $A_N$, all-class macro accuracy, and the harmonic mean in Eq.~\eqref{eq:harmonic}. Both datasets use equal numbers of base and novel classes, so all-class macro accuracy is $(A_B+A_N)/2$. On the balanced CIFAR test set this also equals sample accuracy; on Pets the macro measure gives each breed equal weight. Stability measures include maximum absolute change in an existing score, change of the maximizing class restricted to base classes, and prediction changes over the expanded vocabulary. Error flips are reported as fractions of all base test images, not conditional fractions of previously correct images.

\subsection{Recognition after five-shot class expansion}
\begin{table*}[t]
\centering
\caption{Five-shot expansion from 40 base classes to 80 classes. All values are accuracy percentages, reported as mean $\pm$ sample standard deviation over three distinct class partitions. All readouts share frozen CLIP features; coefficients are selected using base-only calibration. Base and novel performance jointly specify each operating point.}
\label{tab:expansion}
\small\begin{tabular}{lrrrr}
\toprule
Readout & Base & Novel & All & Harmonic \\
\midrule
CLIP text & 78.32 $\pm$ 0.25 & 76.86 $\pm$ 1.97 & 77.59 $\pm$ 0.86 & 77.57 $\pm$ 0.88 \\
Visual prototypes & 77.53 $\pm$ 2.03 & 60.92 $\pm$ 3.94 & 69.23 $\pm$ 2.09 & 68.17 $\pm$ 2.49 \\
Text + prototypes & 82.47 $\pm$ 0.87 & 77.53 $\pm$ 0.47 & 80.00 $\pm$ 0.56 & 79.92 $\pm$ 0.55 \\
Tip-Adapter (tuned) & 78.38 $\pm$ 0.24 & 77.42 $\pm$ 1.72 & 77.90 $\pm$ 0.83 & 77.89 $\pm$ 0.83 \\
Raw evidence fusion & 78.32 $\pm$ 0.25 & 76.86 $\pm$ 1.97 & 77.59 $\pm$ 0.86 & 77.57 $\pm$ 0.88 \\
Candidate-normalized table & 85.36 $\pm$ 0.66 & 70.83 $\pm$ 4.33 & 78.09 $\pm$ 2.02 & 77.36 $\pm$ 2.48 \\
Reference-anchored table & 86.03 $\pm$ 1.40 & 66.90 $\pm$ 4.23 & 76.46 $\pm$ 1.43 & 75.19 $\pm$ 2.14 \\
\bottomrule
\end{tabular}
\end{table*}

Reference-anchored fusion reaches $\RefBase\pm\RefBaseSD$ base accuracy and $\RefNovel\pm\RefNovelSD$ novel accuracy (Table~\ref{tab:expansion}). Its base accuracy is 7.71 percentage points above text-only classification and 3.56 points above text--prototype fusion, providing a retention-oriented operating point for class expansion.

Text--prototype fusion uses visual evidence for every class and reaches $\SimpleAll\pm\SimpleAllSD$ all-class accuracy and $\SimpleNovel\pm\SimpleNovelSD$ novel accuracy. Candidate-normalized fusion reaches $\DynamicBase\pm\DynamicBaseSD$ base and $\DynamicNovel\pm\DynamicNovelSD$ novel accuracy. These points characterize the recognition trade-off accompanying the reference-anchored score system.

Raw fusion selects the text-only point in all five-shot CIFAR partitions. Both normalization rules activate learned evidence in their selected solutions, demonstrating the role of scale alignment when combining classifier logits with cosine similarities.

\subsection{Fine-grained expansion on Oxford-IIIT Pets}
\begin{table*}[t]
\centering
\caption{Five-shot Oxford-IIIT Pets expansion from 15 to 30 breeds under long-tailed base supervision. Values are macro accuracy percentages, mean $\pm$ sample standard deviation across three distinct breed partitions. Grids and the base-only calibration rule are transferred from the CIFAR experiment without dataset-specific search expansion.}
\label{tab:pets}
\small\begin{tabular}{lrrrr}
\toprule
Readout & Base & Novel & All & Harmonic \\
\midrule
CLIP text & 92.87 $\pm$ 3.37 & 94.05 $\pm$ 1.35 & 93.46 $\pm$ 1.59 & 93.43 $\pm$ 1.63 \\
Visual prototypes & 92.08 $\pm$ 2.61 & 68.30 $\pm$ 4.62 & 80.19 $\pm$ 2.33 & 78.35 $\pm$ 2.97 \\
Text + prototypes & 94.77 $\pm$ 2.91 & 93.91 $\pm$ 1.55 & 94.34 $\pm$ 1.43 & 94.31 $\pm$ 1.44 \\
Tip-Adapter (tuned) & 93.56 $\pm$ 2.45 & 94.07 $\pm$ 1.18 & 93.81 $\pm$ 1.27 & 93.80 $\pm$ 1.28 \\
Raw evidence fusion & 90.13 $\pm$ 8.01 & 91.18 $\pm$ 4.27 & 90.66 $\pm$ 6.13 & 90.63 $\pm$ 6.18 \\
Candidate-normalized table & 96.85 $\pm$ 1.50 & 86.49 $\pm$ 7.64 & 91.67 $\pm$ 4.56 & 91.29 $\pm$ 4.99 \\
Reference-anchored table & 97.18 $\pm$ 2.25 & 87.79 $\pm$ 6.18 & 92.49 $\pm$ 2.83 & 92.14 $\pm$ 3.23 \\
\bottomrule
\end{tabular}
\end{table*}

On fine-grained Pets expansion, reference anchoring reaches $\PetRefBase\pm\PetRefBaseSD$ base and $\PetRefNovel\pm\PetRefNovelSD$ novel accuracy (Table~\ref{tab:pets}). Base accuracy is 2.41 points above text--prototype fusion. Mean base/novel accuracy also increases from $\PetDynamicBase$\%/$\PetDynamicNovel$\% with candidate normalization to $\PetRefBase$\%/$\PetRefNovel$\% with reference anchoring.

Text--prototype fusion reaches $\PetSimpleAll\pm\PetSimpleAllSD$ all-class macro accuracy; the reference-anchored table reaches $\PetRefAll\pm\PetRefAllSD$ with greater retention of base-breed discrimination. Both datasets use the same search grids and base-only calibration rule.

\subsection{Calibration controls and analytic incremental readouts}
\begin{table*}[t]
\centering
\caption{Additional five-shot controls with the same frozen encoder, support images, and base-only calibration. Values are mean $\pm$ sample standard deviation across three partitions. The random-projection readout is a controlled frozen RanPAC-style adaptation described in the text.}
\label{tab:controls}
\small
\begin{minipage}{.49\textwidth}\centering CIFAR-100\\[3pt]\resizebox{\linewidth}{!}{\begin{tabular}{lrrr}
\toprule
Readout & Base & Novel & All \\
\midrule
Text + group offset & 78.93 $\pm$ 2.75 & 75.56 $\pm$ 4.40 & 77.24 $\pm$ 0.89 \\
Anchored text + prototypes & 82.40 $\pm$ 2.02 & 76.19 $\pm$ 3.18 & 79.30 $\pm$ 0.59 \\
Frozen RP-ridge (2000) & 70.18 $\pm$ 0.69 & 52.52 $\pm$ 4.27 & 61.35 $\pm$ 1.83 \\
\bottomrule
\end{tabular}}\end{minipage}\hfill
\begin{minipage}{.49\textwidth}\centering Oxford-IIIT Pets\\[3pt]\resizebox{\linewidth}{!}{\begin{tabular}{lrrr}
\toprule
Readout & Base & Novel & All \\
\midrule
Text + group offset & 95.29 $\pm$ 3.80 & 89.85 $\pm$ 4.86 & 92.57 $\pm$ 2.31 \\
Anchored text + prototypes & 94.58 $\pm$ 3.22 & 90.38 $\pm$ 5.41 & 92.48 $\pm$ 1.10 \\
Frozen RP-ridge (2000) & 93.49 $\pm$ 3.32 & 73.70 $\pm$ 5.13 & 83.60 $\pm$ 1.13 \\
\bottomrule
\end{tabular}}\end{minipage}
\end{table*}

Offset-calibrated controls share the same pseudo-novel episode (Table~\ref{tab:controls}). Text + group offset uses base-referenced text scores and a novel offset. Anchored text + prototypes uses fixed base references for both branches and visual evidence for every class. Harmonic selection searches offsets from $-4$ to $4$ in steps of $0.25$ and prototype weights $\{0,0.25,0.5,1,2,4\}$. Fixed-reference text--prototype fusion also satisfies Proposition~\ref{prop:invariance}; the score invariant applies across availability masks.

Anchored text--prototype fusion attains 82.40\% base and 76.19\% novel accuracy on CIFAR, and 94.58\%/90.38\% on Pets. Offset-calibrated text attains 78.92\%/75.56\% and 95.29\%/89.85\%, respectively. Together with the heterogeneous table, these controls separate group calibration, evidence configuration, and the shared fixed-reference property.

The analytic baseline follows the random-projection phase of RanPAC~\cite{mcdonnell2023ranpac}. A fixed Gaussian matrix maps CLIP features to 2000 ReLU coordinates, and $(G+\lambda I)W=Q$ accumulates Gram matrices and class cross-products from base fitting data and novel support. Base validation macro accuracy selects $\lambda\in\{10^{-8},\ldots,10^8\}$; the pseudo episode selects a novel offset. This frozen adaptation uses base validation in place of the original within-task MSE selection. It attains 70.18\%/52.53\% base/novel accuracy on CIFAR and 93.49\%/73.70\% on Pets. Novel statistics refit the analytic readout; candidate-table expansion appends rows to its fixed readout.

\begin{table*}[t]
\centering
\caption{Validation retention-targeted calibration at five shots. The pseudo-base target is 95\% of the closed reference-table pseudo-base accuracy. Among feasible grid points, selection maximizes pseudo-novel accuracy; infeasible partitions use the point with highest pseudo-base accuracy, then pseudo-novel accuracy. Feasible counts refer to validation partitions. Test values are percentages, mean $\pm$ sample standard deviation; test retention is not constrained.}
\label{tab:retention}
\small\begin{tabular}{llrrrr}
\toprule
Data & Readout & Base & Novel & All & Feasible \\
\midrule
CIFAR & Text + group offset & 81.47 $\pm$ 1.13 & 71.49 $\pm$ 5.03 & 76.48 $\pm$ 1.99 & 1/3 \\
CIFAR & Anchored text + prototypes & 82.84 $\pm$ 1.50 & 77.11 $\pm$ 4.46 & 79.98 $\pm$ 1.61 & 3/3 \\
CIFAR & Candidate-normalized table & 84.69 $\pm$ 1.92 & 72.45 $\pm$ 4.14 & 78.57 $\pm$ 1.54 & 3/3 \\
CIFAR & Reference-anchored table & 84.56 $\pm$ 1.31 & 73.88 $\pm$ 5.01 & 79.22 $\pm$ 2.21 & 3/3 \\
Pets & Text + group offset & 95.29 $\pm$ 3.80 & 89.85 $\pm$ 4.86 & 92.57 $\pm$ 2.31 & 3/3 \\
Pets & Anchored text + prototypes & 94.03 $\pm$ 3.37 & 90.96 $\pm$ 5.91 & 92.50 $\pm$ 1.27 & 3/3 \\
Pets & Candidate-normalized table & 97.25 $\pm$ 1.78 & 85.26 $\pm$ 7.25 & 91.26 $\pm$ 4.40 & 3/3 \\
Pets & Reference-anchored table & 97.18 $\pm$ 2.25 & 87.79 $\pm$ 6.18 & 92.49 $\pm$ 2.83 & 3/3 \\
\bottomrule
\end{tabular}
\end{table*}

Under the common 95\% validation retention target, reference anchoring reaches 84.56\% base and 73.87\% novel test accuracy on CIFAR, compared with 82.84\%/77.11\% for anchored text--prototype fusion (Table~\ref{tab:retention}). Both are feasible on all three validation partitions. On Pets, the corresponding test points are 97.18\%/87.79\% and 94.03\%/90.96\%. Text alone reaches the CIFAR validation target in one partition. Targets constrain validation selection; the table reports the resulting test operating points. Complete 95\% and 99\% results accompany the experimental records.

\subsection{Vocabulary growth and score preservation}
\begin{figure*}[t]
\centering\includegraphics[width=.98\textwidth]{stability.pdf}
\caption{Vocabulary growth with five-shot novel prototypes. Each curve averages three true class partitions; error bars show sample standard deviation. Reference-anchored fusion preserves existing scores and base-only rankings. Expanded predictions can still change when a new candidate wins. Maximum score drift is measured across all test queries and existing base classes.}
\label{fig:stability}
\end{figure*}

The vocabulary test appends nested sets of 10, 20, and 40 novel classes to the 40 CIFAR base classes, using sorted novel IDs and fixed five-shot gates. At each stage, we compare existing scores with the original base decision.

Reference anchoring gives zero measured base-score drift and zero base-only ranking churn in all nine CIFAR partition--vocabulary conditions (Fig.~\ref{fig:stability}(a,b)). Candidate normalization changes existing scores by up to approximately 1.29--1.74 normalized units and changes a fraction of base-only predictions. The Pets test appends 3, 7, and 15 breeds and likewise records zero anchored score drift and ranking churn in all nine conditions.

At 40 appended classes, expanded predictions change on $\RefChurn\pm\RefChurnSD$\% of base images, including $\RefErrorFlip\pm\RefErrorFlipSD$\% correct-to-wrong changes (Fig.~\ref{fig:stability}(c)). Wrong-to-correct changes are zero, as specified by Corollary~\ref{cor:competition}. The preserved base scores isolate these changes as competition from appended rows.

Five-session expansion adds eight CIFAR classes or three Pets breeds per session, with two random arrival orders per partition. Across 30 session records per dataset, comparisons of all previously admitted rows, including novel rows, give maximum floating-point score differences of $2.85\times10^{-5}$ on CIFAR and $3.58\times10^{-6}$ on Pets. Both arrival orders produce identical final predictions in every partition under class-ID tie breaking.

\subsection{Support budget and evidence availability}
\begin{table}[t]
\centering
\caption{Reference-anchored expansion across novel support budgets. Accuracy percentages are aggregated across partitions after averaging support subsets within each partition.}
\label{tab:shots}
\resizebox{\linewidth}{!}{\begin{tabular}{lrrrr}
\toprule
Shots & Base & Novel & All & Harmonic \\
\midrule
1 & 84.57 $\pm$ 0.88 & 70.86 $\pm$ 5.41 & 77.72 $\pm$ 2.27 & 77.00 $\pm$ 2.85 \\
2 & 85.65 $\pm$ 1.45 & 70.05 $\pm$ 1.13 & 77.85 $\pm$ 0.67 & 77.04 $\pm$ 0.67 \\
5 & 86.03 $\pm$ 1.40 & 66.90 $\pm$ 4.23 & 76.46 $\pm$ 1.43 & 75.19 $\pm$ 2.14 \\
\bottomrule
\end{tabular}}
\end{table}
\begin{figure*}[t]
\centering\includegraphics[width=.76\textwidth]{shots.pdf}
\caption{Base and novel accuracy at one, two, and five support images per novel class. The same fixed-reference scoring rule is used at every budget, with coefficients selected on the corresponding base-only pseudo episode. Error bars represent variation between class partitions.}
\label{fig:shots}
\end{figure*}

At one, two, and five shots, base accuracy remains between 84.57\% and 86.03\%, while novel accuracy ranges from 66.90\% to 70.86\% (Table~\ref{tab:shots}; Fig.~\ref{fig:shots}). Calibration includes zero branch coefficients, allowing the selected mixture to adapt to the support budget. Each budget defines a separately calibrated operating point: support determines the prototype evidence, and fusion determines its competition with learned base discrimination.

\subsection{Mechanism ablation and operating trade-off}
\begin{table*}[t]
\centering
\caption{Evidence removal at the selected five-shot reference-anchored operating point, with all remaining coefficients fixed. Accuracy percentages are mean $\pm$ sample standard deviation over three class partitions.}
\label{tab:ablation}
\small\begin{tabular}{lrrrr}
\toprule
Readout & Base & Novel & All & Harmonic \\
\midrule
Reference-anchored table & 86.03 $\pm$ 1.40 & 66.90 $\pm$ 4.23 & 76.46 $\pm$ 1.43 & 75.19 $\pm$ 2.14 \\
Without learned evidence & 58.10 $\pm$ 32.21 & 80.72 $\pm$ 4.47 & 69.41 $\pm$ 13.87 & 62.85 $\pm$ 25.08 \\
Without novel prototypes & 86.42 $\pm$ 1.71 & 47.17 $\pm$ 30.85 & 66.79 $\pm$ 14.62 & 56.15 $\pm$ 30.05 \\
Without offset & 86.14 $\pm$ 1.69 & 65.58 $\pm$ 7.35 & 75.86 $\pm$ 2.83 & 74.25 $\pm$ 4.08 \\
\bottomrule
\end{tabular}
\end{table*}
\begin{figure}[t]
\centering\includegraphics[width=.98\linewidth]{tradeoff.pdf}
\caption{Base--novel operating curves for five-shot expansion. The offset is swept with the other selected coefficients fixed; these test curves are descriptive analyses and are not used to select the reported operating points. Symbols identify base-only calibrated points and matched baselines.}
\label{fig:tradeoff}
\end{figure}

Removing learned evidence, novel prototypes, or the offset changes a different part of the selected decision rule (Table~\ref{tab:ablation}). With the remaining coefficients fixed, these interventions identify how each source contributes to the five-shot operating point. The independently calibrated text and text--prototype readouts provide the corresponding reduced-model comparisons.

Figure~\ref{fig:tradeoff} sweeps $d$ from $-2$ to $2$ in steps of 0.5 with $a$ and $u$ fixed. Increasing $d$ shifts the margin between the best base and novel candidates while preserving within-group rankings (Eq.~\eqref{eq:base-novel}). The resulting curve identifies how the offset controls admission of new categories; a validation retention target selects the operating point.

\subsection{Readout efficiency}
Candidate scoring from cached 768-dimensional features takes approximately \LatencyFifty, \LatencySixty, and \LatencyEighty\ ms per batch for 50, 60, and 80 candidates. Measurements use batches of 256 on an NVIDIA RTX 4070 Laptop GPU, five warm-up iterations, and 50 timed iterations with GPU synchronization. They include reference statistics and fusion and exclude image/text encoding. At these vocabulary sizes, launch and reduction overheads dominate the measured latency.

\FloatBarrier
\section{Discussion}
A fixed reference frame makes heterogeneous class evidence reusable as the vocabulary grows. Base-trained columns preserve established discrimination, while semantic anchors and support prototypes score arriving categories in the same coordinate system. The table exposes both the evidence available to a class and the scale on which it contributes.

This separation supports predictable expansion: previously stored rows retain their scores, and new rows enter through candidate competition. Fixed-reference text--prototype fusion shares the invariant, demonstrating its applicability across evidence configurations. Learned base columns supply a retention-oriented operating point within this shared score system.

Base-only calibration makes the recognition objective explicit before novel-class evaluation. Harmonic selection balances the groups, and retention-targeted selection prioritizes a specified pseudo-base accuracy. The vocabulary and sequential tests measure score consistency alongside the recognition outcome at each operating point.

\section{Conclusion}
Reference-anchored candidate tables combine semantic, trained, and support-derived evidence as class inventories expand. Fixed base references decouple normalization from the active vocabulary and preserve existing scores and rankings under append-only expansion. Two-dataset evaluations with base-only calibration verify this stability and characterize retention-oriented recognition across support budgets and sequential arrivals. The resulting interface admits new categories while keeping the existing score system consistent.

\section*{Data and code availability}
The experiments use public CIFAR-100 and Oxford-IIIT Pets data and the publicly released CLIP ViT-L/14 checkpoint. Code, source-hashed experimental records, prediction archives, manuscript sources, and table/figure generation scripts are available at \url{https://github.com/Cheerssom/Reference-Anchored-Candidate-Tables}. Dataset images, pretrained weights, and feature caches are obtained separately using the repository instructions.

\section*{Acknowledgements}
This research received no external funding.
\section*{Competing interests}
The author declares no competing interests.

\bibliographystyle{fcs}
\bibliography{references}

\appendix
\section{Reproduction details}
\label{app:reproduction}
The expansion experiment is implemented in \nolinkurl{src/experiments/fcs_expansion.py}. It can be run from any working directory. The manifest records the code hash, frozen-cache hash, image-pool hashes, class assignments, pseudo-class assignments, support indices, regularization choices, and selection grids. The frozen support features are computed with the same deterministic interpolation used for the frozen training cache. The cached image embeddings are checked numerically before scoring.

Every novel support image comes from the five-image-per-class reserve created before long-tail sampling. This supplies a consistent support budget even for tail classes whose fitting pool contains fewer than five images. At five shots the three subset seeds select the same images in different orders, so the reported standard deviation is calculated across the three distinct class partitions. One- and two-shot results average the three subsets within each partition before computing the partition-level variation.

The main result JSON contains all compared readouts, including evidence-removal ablations. The vocabulary records contain the maximum score drift and both restricted and expanded prediction churn. The operating-curve records contain every swept offset, with no selection based on test outcomes. The efficiency records specify batch size and measured vocabulary size. The manuscript tables are generated directly from these files, and the summary CSV exposes both means and sample standard deviations.

The semantic prompts are ``a photo of a \{\}.'', ``a bad photo of a \{\}.'', ``a cropped photo of a \{\}.'', ``a close-up photo of a \{\}.'', ``a bright photo of a \{\}.'', ``a good photo of a \{\}.'', ``a photo of the \{\}.'', and a second occurrence of ``a photo of a \{\}.''. The repeated prompt follows the existing cached anchor construction and is retained consistently for all methods.

Frozen CLIP features are the sole image representation in the expansion experiments. Six fitting images are independently re-embedded and checked for cosine agreement above 0.999 with the cached features. Image hashes verify that fitting, validation, and reserved-support pools are disjoint. The fixed base text block is computed separately from novel text columns; sequential prototype products can exhibit the reported floating-point residuals when matrix widths change.

\section{Interpreting group competition}
Let $M_B(x)=\max_{b\in\mathcal B}[\widehat q_T(x,b)+a\widehat q_L(x,b)]$ and $M_N(x)=\max_{n\in\mathcal N}[\widehat q_T(x,n)+u\widehat q_P(x,n)]$. The classifier selects a novel class precisely when $M_N(x)+d>M_B(x)$, subject to the fixed tie-breaking rule. Hence, the offset is a threshold on the margin $M_B-M_N$. Increasing $d$ cannot change the winning base class or the winning novel class when all other quantities are fixed. It changes which group supplies the final winner. This explains the direction of the operating curves without assuming a calibrated probability interpretation for the fused score.

For nested candidate sets and fixed references, adding novel rows can increase $M_N$ but cannot change $M_B$. An existing-class example may therefore cross the admission boundary as the vocabulary grows, while its base scores remain identical. This distinguishes the deployment behavior measured by expanded prediction churn from the normalization behavior measured by base-only ranking churn.

\section{Adapted-backbone closed-set configuration}
The candidate-table interface also supports a separate adapted-backbone closed-set configuration. The existing CIFAR-100-LT experiment combines a LoRA-adapted IN21k ViT-B/16 readout with CLIP text and CLIP linear evidence, averages five deterministic test views, and selects gates on held-out training validation data. Three independent adapter/probe training seeds yield 88.86\%, 89.20\%, and 89.55\%, summarized in the original record as $89.20\pm0.28$\% using population standard deviation. This is a 100-class supervised recognition result with different backbone and inference conditions. It illustrates the interface's ability to incorporate adapted evidence; the frozen base-only experiments above establish the class-expansion results.

\end{document}
